% Драфт: Автоматизированный контроль качества компонент кальциевого имиджинга
% Источник: autoinspect_report_ru.docx
% Куда: после «Обработка данных кальциевого имиджинга» в update_nof.tex

\subsection{Автоматизированный контроль качества выделенных компонент}\label{sec:autoinspect}

Как отмечалось выше, после декомпозиции видеозаписей алгоритмом CNMF-E полученные компоненты необходимо классифицировать на нейроны и артефакты. Стандартная процедура предполагает ручную верификацию каждого кандидата экспертом, что требует значительных временных затрат и приводит к межэкспертной вариабельности оценок. Для автоматизации этого этапа была разработана система машинного обучения, обучающаяся на решениях эксперта и воспроизводящая их с высокой точностью.

\paragraph{Пространство признаков.}
Для каждого кандидата в нейроны вычисляется 35 признаков, охватывающих пять категорий:
\begin{itemize}
    \itemsep0em
    \item \textit{Пространственные и морфологические признаки} (11): площадь, выпуклость, соотношение сторон, компактность, эксцентриситет и другие характеристики пространственного следа.
    \item \textit{Временные и статистические признаки} (9): базовая флуоресценция, дисперсия, коэффициенты асимметрии и эксцесса сигнала $dF/F$, автокорреляционные характеристики.
    \item \textit{Событийные признаки} (8): частота кальциевых событий, событийное отношение сигнал/шум, качество аппроксимации при моделировании кинетики кальция, времена нарастания и спада.
    \item \textit{Метрики качества реконструкции} (4): оценка соответствия наблюдаемого сигнала генеративной модели кальциевой динамики.
    \item \textit{Алгоритмические метрики CaImAn} (2): встроенные оценки SNR и пространственной корреляции.
\end{itemize}

Детекция кальциевых событий, необходимая для вычисления событийных признаков, выполняется вейвлет-методом из пакета DRIADA (см.\ раздел~\ref{sec:driada}). Для компонент со сложно характеризуемыми паттернами активности реализована иерархическая система из пяти уровней детекции с последовательным ослаблением критериев качества, обеспечивающая покрытие свыше 99\% кандидатов.

\paragraph{Классификатор.}
В качестве модели классификации использовались объяснимые бустинг-машины (Explainable Boosting Machines, EBM)~\cite{InterpretML-EBM}~--- обобщённые аддитивные модели с парными взаимодействиями~\cite{Lou2013-GA2M}, в которых вклад каждого признака представлен индивидуальной функцией формы:
\begin{equation}\label{eq:ebm}
    g(y) = \beta_0 + \sum_{j=1}^{p} f_j(x_j) + \sum_{(i,j)} f_{ij}(x_i, x_j),
\end{equation}
где $f_j$~--- одномерные функции формы, а $f_{ij}$~--- парные взаимодействия. Такая структура обеспечивает интерпретируемость модели: для каждой классификации можно определить вклад каждого признака в итоговое решение. При этом производительность EBM сопоставима с ансамблевыми методами (случайные леса, градиентный бустинг).

Анализ важности признаков показал, что наибольшую дискриминационную способность обеспечивают временная статистика сигналов (дисперсия, эксцесс) и событийные характеристики (событийный SNR), тогда как встроенные метрики CaImAn, традиционно используемые для фильтрации, оказались менее информативными.

\paragraph{Валидация.}
Система оценивалась методом перекрёстной валидации на наборе из 92\,242 размеченных компонент (187 сессий записи, 4 эксперимента). С учётом асимметрии ошибок классификации (ложноположительные результаты~--- сохранение артефактов~--- более критичны, чем пропуск нейронов) порог принятия решения оптимизировался по метрике $F_\beta$ с $\beta = \sqrt{1/3}$, придающей повышенный вес точности. Достигнутая производительность: точность (precision) 97{,}6\%, полнота (recall) 93{,}3\%, общая метрика $F_\beta = 96{,}5\%$.

Автоматизированная система интегрирована в конвейер BEARMIND и позволяет сократить время контроля качества с нескольких часов до нескольких минут на сессию, обеспечивая при этом воспроизводимость решений.
